% Generated by scripts/assemble.py; corrections belong in source/corrections.json.
\examhead{192}{西安电子科技大学 2026 年数学分析试卷}{181}
% p182-b21; source PDF page 182
\par\Needspace{5\baselineskip}
\blockanchor{p182-b21}
\majorquestion{一}计算题. 每题 7 分, 共 42 分.\par
% p182-b22; source PDF page 182
\blockanchor{p182-b22}
\question{1}设  $0 < x_{n} < 2025, n = 1, 2, \cdots$ ，求  $\lim_{n \to \infty} \frac{x_{n}}{x_{1} + x_{2} + \cdots + x_{n}}$ .\par
% p182-b23; source PDF page 182
\blockanchor{p182-b23}
\question{2}求 $\lim_{x\to 0}\frac{1}{x}\int_{x}^{2x}\sin \frac{1}{t}\mathrm{d}t.$\par
% p182-b24; source PDF page 182
\blockanchor{p182-b24}
\question{3}计算 $\iint_{D} |x^{2} + y^{2} - 2y| \, \mathrm{d}x \, \mathrm{d}y$, 其中 $D: x^{2} + y^{2} \leq 4$.\par
% p183-b00; source PDF page 183
\blockanchor{p183-b00}
\question{4}求 $\iint_{\Sigma} x^{3} \mathrm{d}y \mathrm{d}z + y^{3} \mathrm{d}z \mathrm{d}x + (z^{3} + 3) \mathrm{d}x \mathrm{d}y,$ 其中 $\Sigma: x^{2} + y^{2} + z^{2} = 1, z \geq 0,$ 取下侧.\par
% p183-b01; source PDF page 183
\blockanchor{p183-b01}
\question{5}设  $I_{n}=\int_{0}^{\frac{\pi}{4}}\sin^{n}x\cos x dx,\ n=0,1,2,\cdots$ ，求  $\sum_{n=0}^{\infty}I_{n}.$\par
% p183-b02; source PDF page 183
\blockanchor{p183-b02}
\question{6}设  $f(x)=\left\{\begin{aligned}x,&x\leq0,\\ \frac{1}{n},&\frac{1}{n+1}<x\leq\frac{1}{n},n=1,2,\cdots.\end{aligned}\right.$  求  $f'(0)$ .\par
% p183-b03; source PDF page 183
\par\Needspace{5\baselineskip}
\blockanchor{p183-b03}
\majorquestion{二}判断题. 正确说明理由, 错误给出反例. 每题 6 分, 共 24 分.\par
% p183-b04; source PDF page 183
\blockanchor{p183-b04}
\question{1}若 $a_{n} > 0$ 且 $\varlimsup_{n\to \infty}\frac{a_{n + 1}}{a_n} >1,$ 则级数 $\sum_{n = 1}^{\infty}a_n$ 发散.\par
% p183-b05; source PDF page 183
\blockanchor{p183-b05}
\question{2}任意数列都有单调子列.\par
% p183-b06; source PDF page 183
\blockanchor{p183-b06}
\question{3}设二元函数 $f(x, y)$ 具有连续的偏导数，且 $xf_x + yf_y = 0$，则 $f(x, y)$ 为常值函数.\par
% p183-b07; source PDF page 183
\blockanchor{p183-b07}
\question{4}若函数 $f$ 在任一闭区间 $[a, b]$ 上能取到介于 $f(a)$ 与 $f(b)$ 之间的每一个值，则 $f$ 在 $\mathbb{R}$ 上连续.\par
% p183-b08; source PDF page 183
\par\Needspace{5\baselineskip}
\blockanchor{p183-b08}
\majorquestion{三}(14 分) 证明: $f(x) = \begin{cases} x \left(2 + \sin \frac{1}{x}\right), & x \neq 0, \\ 0, & x = 0. \end{cases}$ 在 $\mathbb{R}$ 上一致连续.\par
% p183-b09; source PDF page 183
\par\Needspace{5\baselineskip}
\blockanchor{p183-b09}
\majorquestion{四}(14 分) 设函数 $f(x)$ 在 $(a, +\infty)$ 内二阶可导, 且 $\lim_{x \to a^{+}} f(x) = \lim_{x \to +\infty} f(x) = A$.\par
% p183-b10; source PDF page 183
\blockanchor{p183-b10}
(1) 若 $A$ 有限, 证明: 存在 $\xi \in (a, +\infty)$, 使得 $f''(\xi) = 0$.\par
% p183-b11; source PDF page 183
\blockanchor{p183-b11}
(2) 若 $A = +\infty$ 或 $-\infty$, 请问 (1) 中的结论是否仍然成立?\par
% p183-b12; source PDF page 183
\par\Needspace{5\baselineskip}
\par\Needspace{6\baselineskip}
\blockanchor{p183-b12}
\majorquestion{五}(14 分) 设 $f(x)$ 连续, 证明:\par
\nopagebreak[4]
% p183-b13; source PDF page 183
\blockanchor{p183-b13}
\displayformula{\iint_ {\Sigma} f (a x + b y + c z) \mathrm{d} S = 2 \pi \int_ {- 1} ^ {1} f (\sqrt {a ^ {2} + b ^ {2} + c ^ {2}} t) \mathrm{d} t.}
% p183-b14; source PDF page 183
\blockanchor{p183-b14}
其中 $\Sigma: x^{2} + y^{2} + z^{2} = 1$.\par
% p183-b15; source PDF page 183
\par\Needspace{5\baselineskip}
\blockanchor{p183-b15}
\majorquestion{六}(14 分) 设 $\alpha > 0$, 正项级数 $\sum_{k=1}^{\infty} a_k$ 收敛, $b_n = \prod_{k=1}^{n} \left(1 + a_k^{1-k^{-\alpha}}\right)$, 证明: 数列 $\{b_n\}$ 收敛.\par
% p183-b16; source PDF page 183
\par\Needspace{5\baselineskip}
\blockanchor{p183-b16}
\majorquestion{七}(14 分) 设 $f(x) = \sum_{n=1}^{\infty} \frac{\cos nx}{\sqrt{n^3 + n}}$, 证明:\par
% p183-b17; source PDF page 183
\blockanchor{p183-b17}
(1) $f(x)$ 在 $(- \infty, + \infty)$ 上连续.\par
% p183-b18; source PDF page 183
\blockanchor{p183-b18}
(2) $f(x)$ 在 $(0,2\pi)$ 上可导.\par
% p183-b19; source PDF page 183
\blockanchor{p183-b19}
(3) 记 $F(x) = \int_{0}^{x} f(t) \, \mathrm{d}t$，则 $\frac{\sqrt{2}}{2} - \frac{1}{15} < F\left(\frac{\pi}{2}\right) < \frac{\sqrt{2}}{2}$.\par
% p183-b20; source PDF page 183
\par\Needspace{5\baselineskip}
\blockanchor{p183-b20}
\majorquestion{八}(14 分) 设 $F_{1}(t) = \sqrt{\pi} \int_{0}^{t} e^{-x^{2}} \mathrm{d}x, F_{2}(t) = \int_{0}^{+\infty} \frac{\sin(tx)}{x} e^{-\frac{x^{2}}{4}} \mathrm{d}x,$ 证明: $F_{1}(t) = F_{2}(t), t \in \mathbb{R}$.\par
